\section{The Order of the Iron Crown}

The Iron Crown of Lombardy, dating from the seventh century, is said to contain one of the nails used in the crucifixion of Jesus. The Holy Roman Emperor \textbf{Charlemagne} used it at his coronation ceremony in Lombardy, as did several succeeding emperors. Even Napoleon used the crown when he claimed to be King of Italy. Two attempts to create an Order of the Iron Crown were made, first by Napoleon, then by King Victor Emmanuel; both came to nothing.

Apparently Julius Evola outlined his own ideas about the Order. Evola both acknowledges the inception of ascetic orders following the fall of the Roman Empire, while denying any link to the Lombardian Iron Crown, despite using the same name. (The published translation is incorrect, either accidentally or deliberately to obscure the ancient connection.)

I'll not bother to comment on the eight points as they are generic and innocuous in themselves. Certainly, it is unlikely to ever be implemented as they stand. While Cordelia for Lear has the good sense to stop at that point, another site adds a section about qualifications and structure.

What is curious about this ascetic order, especially since there are provisions for passing on membership to sons, is a comment about an affiliated formation of women who would serve as sexual surrogates for the men. Methinks Evola is trying to show he's ``got game'', but I don't see the appeal. Maybe the dateless computer geeks who seem to be reading Evola these days relish the idea of 50 men sharing the same half dozen or so thots. This would be sure to lead to attachments and the inevitable jealousies and conflicts that would disrupt the group. Any man with real world experience, instead of living in his mother's apartment, can grasp this. Believe me, the women will turn out to be very demanding, to the point of diverting attention from the order's purpose.

Then there is the delicate problem of recruitment. Perhaps in decadent Rome, such an order can find suitable women. But in Evola's ancestral Sicily, those women are likely to be some man's daughter or sister. I can guarantee this idea would not have gone over very well and the Sicilians had their own secret order to deal with such issues.

We have made available several posts about actually existing knightly orders, for example at Random Thoughts on Chivalry\footnote{\url{https://gornahoor.net/?p=277}}. Notice the profound difference in values from what Evola proposes. I suggest that anyone interested in a chivalric or ascetic order should take a look at previously successful endeavors and not the laundry list of a man who was never part of any order.


\flrightit{Posted on 2010-04-28 by Cologero}

\begin{center}* * *\end{center}

\begin{footnotesize}
\begin{sffamily}
\texttt{James O'Meara on 2010-05-03 at 17:57 said:}

I've frequently commented on my blog regarding Evola's rather eccentric ideas on sexuality, so this caught my eye. I don't see this as particularly shocking or unworkable. The footnote makes it clear that this is in the context of sexual alchemy, not recreation or `scoring.' The idea female for this activity, as Evola notes in \textit{Metaphysics of Sex}, would be a virgin suitable for deflowering; as this is a one-time event, the whole issue of ``jealousy'' etc. doesn't arise. I suppose the women would be disposed of in some de Sadean way; this would be consistent with his other preferred activity, flagellation.

``I realize these are unpopular idea'' as a Python character says, but they are not unusual in the alchemical tradition; just as Madonna knows nothing about Qabalah, Sting knows nothing about Tantra. Actual Tantric or Taoist practices involved the ruthless exploitation of women for their ``essence'' as General Jack Ripper would say. Women represent ``power'' but that is precisely why they must be controlled and assimilated; no ``goddess'' worship here.

You may find these ideas untrue i.e., of no practical use or ``immoral'' on the basis of some feminine, Christian morality, but they are hardly ``unusual'' in Tradition.

That a celibate order would need to provide for the children of its members would not be surprising to anyone familiar with the Renaissance Church; for the same reason, I would assume Sicly would in fact be an excellent source of virgins ask Michael Corleone where their ``protective'' family would, for that very reason, be able to hand them over to the local Elites. Happens every day, man.

\hfill

\texttt{Cologero on 2010-05-03 at 23:53 said:}

I guess I just find it unlikely. Evola did write of a third order, so I understand that to mean a more or less permanent underclass of groupies, not disposable virgins. In order to count on the locals supplying the virgins willingly, the Order would have to be doing favours in return ... just ask Michael Corleone. The actual activities of the order are left undefined.

I wasn't addressing the post so much to Evola, but rather to rather to a certain type of neo-pagan who has not a chance at a virgin. For example, consider this quote from a Odinist facebook page: ``The key to happiness is a goblet full of mead, a wench, and a steed.'' If sex and intoxicants were the key to spirituality, there would be pilgrimages to every American ghetto. But this Odinist wants to act like one; just needs his low rider to make him a man.

On the contrary, sex magick is only of value to someone who has no need for sex (as Evola subtly implies). When the Dalai Lama was asked about Tantric initiation (naturally of intense interest to his decadent American followers), he pointed out that it is only appropriate to someone who could dine on a plateful of shit and a glass of pus or a gourmet meal with the same indifference.

Even our friend \textit{Aleister Crowley}\footnote{\url{http://www.gornahoor.net/?p=51}}, in a lucid moment wrote:

\begin{quotex}\footnotesize

It is, indeed, of the first importance for the celebrant in any phallic rite to be able to complete the act without even once allowing a sexual or sensual thought to invade his mind. The mind must be as absolutely detached from one's own body as it is from another person's.
\end{quotex}


Who do you know who is capable of that? He's the one that gets to deflower the Virgin. But that attitude is what indeed makes it ``unusual'' in Tradition, except, perhaps, for poseurs.

But aren't there other ways to deal with feminine energy. Do you suppose \textbf{Moonchild} is a bit too sweet?


Or the idea of Polar beings to bring completion. There may be something necessary in the constant irritant of a female companion. Just hanging out with the guys is an option, but even the misogynist went on the quest for ``She who must be obeyed.'' (Rider Haggard)

\hfill

\texttt{Matt on 2010-05-08 at 20:53 said:}

The difference in Evola's list of requirements and the requirements of the medieval orders is quite interesting. Your suggestion of a possible inaccurate translation is possible, though there is also the possibility that Evola didn't want to include elements that he thought were of a christian nature.

The last sentence of this blog entry makes me think of his ur group and I have always wondered whether Evola and the group at large achieved any kind of inner/spiritual realization, at least of any considerable magnitude.

\hfill

\texttt{Francis Mercuri on 2010-05-10 at 03:38 said:}

Mark, while it is somewhat true that mainstream Christianity enjoys a greater, albeit superficial, acceptance as an organized religion, as a ``norm'', among the general Western public (as opposed to Paganism), this does not imply that traditional Christians are not ``tested''.

1. In the modern (predominantly liberal) social (``public'') domain, which is characterized by materialism, mechanism, moral relativism, and to a large extent atheism, the ideals of traditional Christianity are perceived as antagonistic to the overall operation of the current order. So, it is ``ok'' to be Christian, just as long as one ``keeps their opinions to themselves'', and doesn't try to apply their traditional Christian theology, philosophy, or ethics to ``real world'' scenarios, that might just ``upset the apple cart''. As far as modernity is concerned, the world of Medieval Catholicism, or Eastern Orthodoxy, is just as anathema as ``Paganism'', if not more so.

2. Inwardly, the authentic Christian path is entirely about being ``tested''. I get the feeling that you understand Christianity as being just a ``belief'', or a sentimental ``feeling'', sort of akin to how one might acquiesce to popular sentiments about their home team(s). However, if you dig a bit deeper, you will find that Christianity is about winning a precious medicine. This healing balm repairs the rupture between a divided immanent and transcendent in man, the fundamental cause originating with the ``Fall'', but the secondary and ongoing causes arising in man himself, as the result of attachment (identification with) the passions, or as some might say, the innumerable ``i's'', which obscure the true ``Personality'' (per Guenon's meaning of ``Personality''). These forces, are what the Christian contends with, in an ongoing effort of ``spiritual athletics''--not in the sense of being ``sins'' or violated rules, to be punished (as exoteric jargon might phrase it), but as energies obfuscating the ``Real Presence''. This really differs little from what Evola writes about how to deal with the ``emotions'' while pursuing the real ``I'', in his tract about the Mithraic Mysteries.

In the past, I've explained things this way to Pagans, who marginally agree with the content as expressed, but always with the retort ``Yes, but, that isn't how the `average' Christian understands Christianity'', and I anticipate this is what you might be thinking. In response, I'd say that simply it matters little what the `average' Christian thinks in this regard, as what we are talking about is Christianity ``as it is''; that is to say, not what some modern Protestant deviation has made of it--but, what it says of itself, via the totality of Scripture, the Orthodox Liturgies/Liturgical cycle, Patristic and Saintly writings/exegesis, ascetical commentaries such as Philokalia and the Gerontikons, the symbolism inherent in Iconography, among other sources.

Minimally, in these ways, I do consider that the Christian is indeed ``tested''; but, like the owner of the forum recently stated, we must go beyond the ``superficial distinctions'' to grasp this. It is a waste of energy on all parties behaves to beat each other over the head concerning false dichotomies surrounding some endless ``war'' between a traditional Christianity and Paganism.

\hfill

\texttt{Cologero on 2010-05-13 at 13:31 said:}

To quote Guenon:

\begin{quotex}\footnotesize

A still more surprising thing is that some people in these days think that they can expound traditional doctrines by taking profane instruction itself as a sort of model, without taking the least account of the nature of traditional doctrines and of the essential differences which exist between them and everything that is today called by the names of ``science'' and ``philosophy'', from which they are separated by a real abyss in so doing they must of necessity distort these doctrines completely by over-simplification and by only allowing the most superficial meaning to appear, otherwise their pretensions must remain completely unjustified.
\end{quotex}


Similarly, it is quite surprising to me, also, that anyone would take notice of this blog while at the same time rejecting out of hand any of its content.

The regurgitation of definitions purloined from a philosophy textbook is passed off as intelligence, and basing one's life someone else's conception published in a pop pychology magazine is considered creativity.

That, sad to say, is the state of the world today. What we have to offer on this blog is a much deeper and richer conception of life. As such, however, it will be accessible only to a superior type of man. They will know themselves by their almost magnetic attraction to traditional doctrines, even if their understanding is at yet immature. The purpose of this blog is twofold:

(1) to counteract the misconceptions aroung Tradition

(2) to encourage those interested to pursue Tradition in depth

The poseurs will eventually get bored and leave.

\hfill

\end{sffamily}
\end{footnotesize}

